\subsection{Coverage and Retention}
\label{app:retention-definition}
Retrieval coverage and retention measure different properties. Recall@50 measures how much of the required skill set is available in the ranked candidate pool. Retention measures how much of that coverage remains when only the first $k$ skills are used. For a method $m$ and retriever $r$, let $\overline{\operatorname{Recall}}_{m,r}@k$ denote recall averaged over test tasks. The retention curves use
\[
\operatorname{Ret}_{m,r}@k
=100\frac{\overline{\operatorname{Recall}}_{m,r}@k}
{\overline{\operatorname{Recall}}_{m,r}@50}.
\]
This is a ratio of task-averaged recalls, rather than an average of task-level retention ratios. We evaluate $k\in\{1,5,10,25,50\}$, and average the three retriever-specific retention values to obtain the curve in Figure~\ref{fig:average_recall_degradation}. Each method is normalized by its own Recall@50. In particular, SkillDAG expands its initial retrieval set, so its normalization uses the expanded ranked pool.

A method can have high retention but low retrieval coverage. Conversely, a method can recover more required skills at cutoff 50 but rank them too low for a small skill budget. We therefore report Recall@50 alongside retention, rather than interpreting the normalized curve alone as end-to-end recall.
